\documentclass[a4paper,11pt]{article}

% Font and encoding
\usepackage{mathptmx}
\usepackage[utf8]{inputenc}
\usepackage[T1]{fontenc}

% Page geometry
\usepackage[margin=2cm,top=2.5cm,bottom=2.5cm]{geometry}

% Compact spacing
\usepackage{setspace}
\setstretch{1.05}
\usepackage{parskip}
\setlength{\parskip}{3pt}
\setlength{\parindent}{0pt}

% Section formatting
\usepackage{titlesec}
\titleformat{\section}{\large\bfseries}{}{0em}{}[\vspace{-4pt}\rule{\textwidth}{0.5pt}]
\titleformat{\subsection}{\normalsize\bfseries}{}{0em}{}
\titleformat{\subsubsection}{\normalsize\itshape\bfseries}{}{0em}{}
\titlespacing*{\section}{0pt}{12pt}{4pt}
\titlespacing*{\subsection}{0pt}{8pt}{2pt}
\titlespacing*{\subsubsection}{0pt}{6pt}{2pt}

% Lists
\usepackage{enumitem}
\setlist{nosep,leftmargin=1.5em,topsep=2pt}
\setlist[itemize]{label=\textbullet}

% Tables
\usepackage{booktabs}
\usepackage{array}

% Hyperlinks and TOC
\usepackage[hidelinks]{hyperref}

% Math
\usepackage{amsmath}

% Header/footer
\usepackage{fancyhdr}
\pagestyle{fancy}
\fancyhf{}
\fancyhead[L]{\small\textit{ESADE Entrepreneurial Finance --- Study Guide}}
\fancyhead[R]{\small\textit{Sessions 00--04}}
\fancyfoot[C]{\small\thepage}
\renewcommand{\headrulewidth}{0.4pt}

\begin{document}

% ============================================================
% TITLE PAGE
% ============================================================
\begin{titlepage}
\vspace*{4cm}
\begin{center}
{\LARGE\bfseries ESADE Entrepreneurial Finance}\\[0.8cm]
{\Large Study Guide}\\[0.5cm]
{\large Sessions 00--04 + Quiz Prep}\\[2cm]

\begin{tabular}{rl}
\textbf{Course:} & Entrepreneurial Finance \& Venture Capital \\[4pt]
\textbf{Professors:} & Carles Florensa, Teresa Corrales, Benjamin Ullrich \\[4pt]
\textbf{Term:} & Term 2 --- March 2026 \\
\end{tabular}
\end{center}
\vfill
\end{titlepage}

% ============================================================
% TABLE OF CONTENTS
% ============================================================
\tableofcontents
\newpage

% ============================================================
% SESSION 00
% ============================================================
\section{Session 00 --- Introduction (Feb 17)}
\textbf{Instructors:} Teresa Corrales, Carles Florensa \hfill \textbf{Case:} Too Good To Go

\subsection{Financial Statements Fundamentals}

\textbf{Balance Sheet} --- financial position at a moment in time. Assets = Liabilities + Owner's Equity.

\textbf{Income Statement (P\&L)} --- tracks revenue and expenses over a period. Net profit flows to equity on the balance sheet.

\textbf{Cash Flow Statement} --- tracks every cash movement in/out; reconciles with balance sheet.

\textbf{Key principle:} Profit does NOT equal Cash --- timing differences (AR, AP) and working capital changes create the gap.

\textbf{CapEx vs OpEx} --- CapEx = investment in assets with $>$1 year lifespan, depreciated over useful life. OpEx = recurring expenses, hit P\&L immediately.

\textbf{Working Capital} = Inventories + Accounts Receivable $-$ Accounts Payable --- an increase = financial need (reduces operating cash flow).

\textbf{Depreciation} --- reduces fixed asset value and earnings over time, but does NOT affect cash.

\textbf{Accrual vs Cash accounting} --- 100\% legal to show accrual to investors and cash for tax filing.

\textbf{FIFO vs LIFO} --- LIFO reduces taxable income; at scale the difference is significant.

\subsection{Core Formulas}

\begin{tabular}{@{}ll@{}}
\toprule
\textbf{Formula} & \textbf{Definition} \\
\midrule
Assets & $=$ Liabilities $+$ Equity \\
COGS & $=$ Units Sold $\times$ Unit Cost \\
Gross Profit & $=$ Sales $-$ COGS \\
Net Profit & $=$ Gross Profit $-$ Expenses \\
Ending Inventory & $=$ Beginning $+$ Purchases $-$ COGS \\
Working Capital & $=$ Inventories $+$ AR $-$ AP \\
\bottomrule
\end{tabular}

\subsection{Financial Modeling (Too Good To Go Case)}

Four pillars:
\begin{enumerate}
\item Business model / CapEx needs
\item Operations --- revenues, expenses, margins
\item Growth projections
\item Working capital policies
\end{enumerate}

\textbf{Key learnings:} Always use formulas, never hardcode. Sensitivity analysis identifies critical variables. Supplier payment terms have massive impact on cash flow. Ending cash must always be positive.

\subsection{Assessment}

\begin{tabular}{@{}lr@{}}
\toprule
\textbf{Component} & \textbf{Weight} \\
\midrule
Class Participation & 30\% \\
Group Assignments (Midterm 12\% + Final 28\%) & 40\% \\
Individual Test (MCQ, Session 9) & 30\% \\
\bottomrule
\end{tabular}

\newpage

% ============================================================
% SESSION 01
% ============================================================
\section{Session 01 --- VC, PE \& Business Angels (Mar 3)}
\textbf{Instructor:} Carles Florensa \hfill \textbf{Case:} AB-Biotics

\subsection{Venture Capital Fundamentals}

VC main objective: Make profit by selling stake in medium term (NOT lending money or running companies).

\begin{tabular}{@{}ll@{}}
\toprule
\textbf{Metric} & \textbf{Value} \\
\midrule
Management fee & 2--2.5\% of fund size \\
Expected return per investment & 40\% IRR over 7--8 years ($\approx$10x multiple) \\
Average fund return & 13--15\% IRR \\
Hurdle rate for LPs & 6--10\% \\
\bottomrule
\end{tabular}

\vspace{4pt}
\textbf{Business angels vs VCs:} Angels invest own money, often for enjoyment and learning, not just returns.

\textbf{Convertible notes:} Debt that converts to equity later. Includes cap (max valuation) and discount. Delays valuation decision until milestones. Faster/simpler than equity rounds.

\subsection{VC vs PE vs CVC}

\begin{tabular}{@{}p{0.15\textwidth}p{0.75\textwidth}@{}}
\toprule
\textbf{Type} & \textbf{Characteristics} \\
\midrule
VC & Early-stage, minority stakes, high-growth tech/innovation, targeting 10x+ returns \\
PE & Later-stage/buyouts, majority control, mature businesses, targeting 2--3x returns \\
CVC & Corporate venture arms, strategic alignment with parent company, may accept lower financial returns for strategic value, risk of conflicts of interest \\
\bottomrule
\end{tabular}

\subsection{VC Fund Structure \& Lifecycle}

\begin{itemize}
\item \textbf{LPs} (limited partners) provide capital; \textbf{GPs} (general partners) manage the fund
\item Management company handles operations
\item \textbf{Capital calls:} LPs commit but money is drawn down over time
\item 5-year investment period + 5-year divestment period = \textbf{10-year fund life}
\item Management fees (2--2.5\%) reduce investable capital --- a \texteuro100M fund may only deploy \texteuro80M
\item \textbf{Carry} (20\% of profits above hurdle) is the GP's real upside
\item GP typically invests 1--2\% as ``skin in the game''
\end{itemize}

\subsection{Business Angel Framework}

\begin{itemize}
\item \textbf{3Fs:} Friends, Family, Fools --- earliest non-institutional capital
\item \textbf{Valley of Death:} gap between initial funding and revenue generation where most startups fail
\item \textbf{Smart Money:} capital that comes with expertise, network, and operational support
\item \textbf{Angel types:} active, latent, virgin, founder angel, super angel
\item \textbf{Harvest events:} IPO, M\&A, secondary sale
\end{itemize}

\subsection{Venture-Stage Financing Map}

\begin{tabular}{@{}ll@{}}
\toprule
\textbf{Stage} & \textbf{Typical Sources} \\
\midrule
Pre-seed / Seed & Grants, 3Fs, angels, accelerators \\
Startup & Angels, seed VCs, convertible notes \\
Growth & Series A/B VCs, venture debt \\
Maturity & PE, mezzanine, bank loans, factoring \\
Exit & IPO, M\&A, secondary \\
\bottomrule
\end{tabular}

\subsection{Fundraising Strategies}

\begin{itemize}
\item \textbf{Bootstrapping:} Maximum control, minimal dilution, but risk of irrelevance if competitors raise big.
\item \textbf{Seed strapping:} Single round to reach profitability, limited dilution, ideal for capital-efficient models.
\item \textbf{Traditional VC:} Necessary when competitors raise significant funds.
\end{itemize}

\subsection{AI and Market Dynamics}

Europe: 35\% of VC to AI; US: 65\%. Traditional later-stage VCs moving to seed/pre-seed. Creates upward pressure on valuations.

\subsection{Case: AB-Biotics}

Founded 2008, needed EUR 1M. Chose business angel (Lafuente) over VCs because of higher valuation (EUR 4M pre-money vs EUR 2M). Exit 2019: acquired by Kaneka for EUR 64M (14.4x multiple). Founders retained 6\% with full autonomy.

\textbf{Decision framework:} Angel (Lafuente) offered biotech expertise, easier governance, less exit pressure. VCs offered signaling, more capital, fundraising/exit support. Angel followed with another EUR 1M; company raised EUR 4.3M at IPO in 2010.

\subsection{Guest Speaker: Ramiro (Keon)}

Founded Keon 15 years ago --- intelligent retail with RFID. EUR 31M revenue, 96\% exports. Sold for EUR 65M (14.4x, 20x with earnout).

\textbf{Key tips:}
\begin{itemize}
\item Delay first round as long as possible. Use public grants with zero interest.
\item Have EUR 200K in sales before first round. Start fundraising 12 months early.
\item Never overvalue early --- creates down round problems.
\item Investors invest in people. Be transparent with financials.
\item Venture debt available near positive EBITDA (9--10\%, $\sim$1\% equity).
\item For exit: hire M\&A banker (3\% fee); industrial buyers pay more.
\item Hardware + software = stronger moat but requires more capital (prototypes, manufacturing, warehouse, working capital).
\item Investor selection should match company's time horizon; angels/family offices may fit better than VC/PE for longer-build businesses.
\item Transparency with investors, decision-focused board meetings.
\item Vendor due diligence before exit; run a competitive M\&A process.
\end{itemize}

\newpage

% ============================================================
% SESSION 02
% ============================================================
\section{Session 02 --- Deal Flow \& Screening (Mar 6)}
\textbf{Instructor:} Carles Florensa \hfill \textbf{Case:} Skilfill

\subsection{Deal Flow Framework}

Sourcing $\rightarrow$ Screening $\rightarrow$ Partner Review $\rightarrow$ Due Diligence $\rightarrow$ Investment Committee $\rightarrow$ Capital Deployment

\textbf{Screening criteria:} team, market, product, business model, traction, fund fit.

\subsection{Pitch Framework}

A strong pitch must: (1) signal legitimacy, (2) reduce uncertainty, (3) tell a compelling narrative.

\subsection{Investment Evaluation}

Investors prioritize \textbf{TEAM} over market, technology, or ideas. Serial entrepreneurs get more support.

\textbf{Four evaluation buckets:} Competence, Complementarity, Psychological traits, Commitment.

\textbf{Biases:} Pedigree, education, pattern matching.

\subsection{Case: Skilfill}

Skills-based hiring AI. Seeking EUR 600K at EUR 4.8M pre-money. Well-structured pitch but critical weaknesses:
\begin{itemize}
\item Inconsistent customer claims (5 in pitch vs 2 in Q\&A)
\item Missing moat
\item Unrealistic LTV/CAC (\$2,887)
\item No infrastructure costs for 2 years
\item No sales budget despite growth targets
\item Revenue based on growth rate not funnel
\item Incomplete P\&L
\end{itemize}
\textbf{Verdict:} Do NOT invest, come back in 6 months with traction.

\textbf{Additional weaknesses:}
\begin{itemize}
\item Pricing inconsistency: pitched EUR 50 vs model EUR 80--96
\item Hidden ``Connect'' product driving large share of projected revenue
\item TestGorilla as competitive threat (scale, funding, customer base) makes moat much thinner
\item \textbf{Investor takeaway:} convertible note with milestones would have fit the risk better than straight equity at EUR 4.8M pre-money
\end{itemize}

\subsection{Guest Speaker: Jose Antonio Gonzalez (Business Angel)}

Deloitte partner 31 years. Best Spanish Business Angel 2025. 85 startups invested, $\sim$10/year at EUR 10--25K. Expects 60+ of 85 to fail. 10--12 year horizon.

\textbf{Philosophy:}
\begin{itemize}
\item Team is critical. Never rush due diligence.
\item Founders must understand numbers. Work 5+ years in industry first.
\item Asks for 5-year org chart. Looks for realistic projections.
\item Must have month-by-month Excel. Don't blackmail family into investing.
\item Know worst case scenario.
\item Pre-seed valuation discipline --- max EUR 2M valuation.
\item Uses 5-year org chart as a founder test --- reveals whether founders think long-term and understand scaling.
\item Validate with customers AND competitors before fundraising.
\end{itemize}

\newpage

% ============================================================
% SESSION 03
% ============================================================
\section{Session 03 --- The Financial Plan (Mar 10)}
\textbf{Instructor:} Carles Florensa \hfill \textbf{Case:} Fuertes Vientos Wind Farm

\subsection{Operating Cash Flow}

\begin{center}
\fbox{\parbox{0.7\textwidth}{\centering
\textbf{OCF = EBIT $-$ Taxes + D\&A $-$ Increase in Working Capital}\\[4pt]
\small If working capital decreases, it is added (not subtracted).
}}
\end{center}

\subsection{Bottom-Up Forecasting Sub-Frameworks}

\begin{itemize}
\item \textbf{CAC-based:} Marketing budget / CAC = acquirable customers
\item \textbf{Funnel/conversion:} Traffic $\rightarrow$ qualified leads $\rightarrow$ opportunities $\rightarrow$ closed deals
\item \textbf{Cohort retention:} Track how many customers from each cohort remain over time
\item \textbf{Cohort-based cash flow:} Different cohorts may have different payment terms
\end{itemize}

\subsection{Revenue Quality Framework}

\begin{itemize}
\item Churn-adjusted growth is more meaningful than topline growth alone
\item 50\% growth with 0\% churn vs 25\% churn = completely different businesses
\item CAC usually rises as firms scale (organic $\rightarrow$ paid channels)
\end{itemize}

\subsection{``Dying of Success'' Problem}

Companies can be profitable but face bankruptcy if working capital grows too quickly. Cash tied up in inventory and AR even as revenue increases.

\subsection{Working Capital Formulas}

\begin{tabular}{@{}ll@{}}
\toprule
\textbf{Item} & \textbf{Formula} \\
\midrule
Inventory & $= (\text{COGS} \times \text{Number of Days}) \;/\; 360$ \\
Accounts Receivable & $= (\text{Sales Revenue} \times \text{Credit Days}) \;/\; 360$ \\
Accounts Payable & $= (\text{COGS} \times \text{Payment Days}) \;/\; 360$ \\
\bottomrule
\end{tabular}

\vspace{4pt}
\textbf{Note:} 360 days is the financial convention.

\subsection{Negative Working Capital (Ideal)}

\textbf{Supermarkets:} Paid immediately, pay suppliers 120+ days.\\
\textbf{SaaS:} No inventory, prepaid subscriptions.\\
Both: working capital finances the business.

\subsection{Free Cash Flow}

The number banks look at for loans. Calculated before financial decisions. If WC increases, FCF decreases.

\subsection{Tax Credit / Loss Carry-Forward}

When companies have losses, they accumulate the right to offset future taxes. Apply this to startup financial models.

\subsection{Balance Sheet Rules}

Assets must ALWAYS equal Equity + Liabilities. Never present unbalanced. Never use a ``hammer'' to force balance.

\subsection{Case: Fuertes Vientos}

Wind energy, operations 2007. 90\% regulated revenue. No inventory. AR 45 days, AP 30 days, Tax 30\%. Depreciation: intangible 5yr, tangible 20yr. Cannot depreciate before operations begin. Loan EUR 30M at 7\%, 10yr linear, 2yr grace period.

\textbf{Operating mechanics:} 42.5 MW $\times$ 2,510 hours capacity. Regulated-rate step-down after year 5. Reinvestment from 2011 onward equal to annual depreciation after warranty expiry.

\textbf{Core strategic takeaway:} Could staged financing have preserved more founder control instead of selling 80\% too early?

\subsection{Guest Speaker: Adria Roca (Inverrere VC)}

Manages EUR 2.2B, VC team $\sim$EUR 200M. Series A focus, EUR 1--5M revenue. Reviewed 17,000 companies, invested in 17 (0.1\%). 50\% of portfolio fails; need one 30x return. 50\% of VC funds don't return capital to LPs.

\textbf{Forecasting:} BIGGEST MISTAKE = top-down approaches. Correct = bottom-up (capacity, conversion rates, sales cycles). CAC is NOT stable --- increases as companies scale. Revenue quality matters more than growth rate.

\textbf{Fundraising:} Avoid premature dilution. Maintain safety buffers. Combine equity + debt. VCs limit dilution to 20--25\%/round. VCs need $>$10\% at exit.

European exits tend to be earlier and lower than US due to smaller markets, limited growth capital, founder fatigue.

\textbf{AI moats:} Increasingly come from proprietary data, embedded enterprise relationships, and compliance/security --- not generic AI features. Growth quality $>$ raw growth, especially when churn and rising CAC are factored in.

\newpage

% ============================================================
% SESSION 04
% ============================================================
\section{Session 04 --- Valuation of Start-ups (Mar 12)}
\textbf{Instructor:} Carles Florensa \hfill \textbf{Cases:} SUN Microsystems, Crypto A \& B

\subsection{VC Method --- The Core Logic}

Start from expected exit value. Work backwards using target return multiple. Result: implied post-money valuation today.

\subsection{Key Definitions}

\begin{tabular}{@{}ll@{}}
\toprule
\textbf{Term} & \textbf{Formula} \\
\midrule
Post-money & $=$ Investment $+$ Pre-money \\
Pre-money & $=$ Post-money $-$ Investment \\
Ownership \% & $=$ Investment $/$ Post-money \\
\bottomrule
\end{tabular}

\subsection{Share-Based Valuation Formulas}

\begin{itemize}
\item Pre-money $=$ pre-existing shares $\times$ new round price per share
\item Post-money $=$ total shares (pre-existing $+$ new) $\times$ price per share
\item \textbf{These are very testable on the exam.}
\end{itemize}

\subsection{5-Step VC Method}

\begin{enumerate}
\item \textbf{Enterprise Value at Exit:} Calculate EV at Year N (e.g., using P/E multiples).
\item \textbf{Investment Value at Exit:} Investment $\times (1 + \text{IRR})^{\text{years}}$
\item \textbf{Required Ownership:} Investment at exit $/$ EV at exit
\item \textbf{Calculate Shares:} $(\text{Fraction} \times \text{pre-existing shares}) \;/\; (1 - \text{Fraction})$
\item \textbf{Price per Share:} Investment $/$ Shares
\end{enumerate}

\subsection{Quick Valuation (for Entrepreneurs)}

\begin{align*}
\text{Post-money} &= \frac{\text{Future Enterprise Value}}{(1 + \text{IRR})^{\text{years}}} \\[4pt]
\text{Pre-money} &= \text{Post-money} - \text{Investment} \\[4pt]
\text{Post-money} &= \text{Investment} \;/\; \text{Dilution\%}
\end{align*}

\subsection{IRR Formula}

\[
\text{IRR} = \left(\frac{\text{Value at exit}}{\text{Investment}}\right)^{1/\text{years}} - 1
\]

\subsection{Price vs Value}

Price can exceed valuation (hype, competing investors). Price can fall below valuation (low liquidity, desperation).

\subsection{Case: SUN Microsystems}

Series A investors: 87.2x multiple at IPO. Founders: $\sim$4,800x. Series H investors (1.5yr pre-IPO): only 14\% IRR.

\textbf{Timing and return dispersion lesson:} founders/early investors made extraordinary returns, late investors earned weak returns. Valuation collapse: \$12.8B (Jan 2008) $\rightarrow$ \$7.4B Oracle acquisition (Apr 2009).

\subsection{Case: Crypto A (No Dilution)}

Sales \$3M, 55\% growth, 30\% margin, 15x P/E. Year 3 EV: \$50.25M. \$2M investment at 40\% IRR = \$5.49M at exit. Ownership: 10.92\%. Shares: 134,845 at \$14.83/share. Pre-money: \$16.31M.

\subsection{Case: Crypto B (With Dilution)}

Accounts for second round. Shares increase, price per share drops to \$13.47.

\textbf{Key lesson:} Future rounds should be modeled upfront because anticipated dilution lowers the max acceptable price/share even if target ownership is unchanged.

\subsection{OpenAI Valuation (Real-World)}

\$110B raised at \$730B pre-money = \$840B post-money. Dilution: 13\%. Revenue multiple: 24x on 2026 forecast of \$30B. AI companies: 20--30x sales. US valuations 2--3x higher than European.

\subsection{Valuation Triangulation}

VC method is one tool, not the only one. Also use: comparable company multiples, comparable transactions, and DCF (when forecasts are credible).

\subsection{Stock Option Pool}

Should dilute ALL shareholders proportionally, not just founders.

\newpage

% ============================================================
% QUIZ PREP
% ============================================================
\section{Quiz Prep --- Practice Questions}

\subsection{VC Investment Process}

\begin{tabular}{@{}p{0.65\textwidth}p{0.3\textwidth}@{}}
\toprule
\textbf{Question} & \textbf{Answer} \\
\midrule
Q1: Final step of VC process & Exit \\
Q2: Key element of investment thesis & Sector and go-to-market approach \\
Q3: NOT part of thesis & Percentage of stake proposed \\
Q4: NOT critical for team evaluation & Money they invested \\
\bottomrule
\end{tabular}

\subsection{Deal Flow}

\begin{tabular}{@{}p{0.65\textwidth}p{0.3\textwidth}@{}}
\toprule
\textbf{Question} & \textbf{Answer} \\
\midrule
Q5: Primary rejection reason & Lack of fit with strategy \\
Q6: Soft screening factor & Trust and culture \\
Q7: Unit economics & Very relevant if you can verify calculations \\
\bottomrule
\end{tabular}

\subsection{Financial Plan}

\begin{tabular}{@{}p{0.65\textwidth}p{0.3\textwidth}@{}}
\toprule
\textbf{Question} & \textbf{Answer} \\
\midrule
Q8: OCF formula & Beginning Cash + EBIT $-$ Taxes + D\&A $-$ Investment in WC \\
Q9: WC Investment (Y1$=$110, Y2$=$140) & Investment $=$ 30 \\
Q10: Days in AP: $(50/360) \times 365$ & 50 days \\
Q11: AR end of year: $(500/365) \times 30$ & EUR 41 \\
Q12: OCF requires\ldots & Taking away the investment in WC \\
Q13: CF impact of AR growth & EUR 275,000 increase \\
Q14: Retained earnings reflected in\ldots & Balance Sheet \\
\bottomrule
\end{tabular}

\subsection{Quiz Convention Note}

Some practice questions use 365-day calculations while the in-class convention uses 360 days. Watch for this in exam wording.

\subsection{Additional Quiz Concepts}

\begin{itemize}
\item \textbf{3Fs, valley of death, smart money} --- exam-friendly definitions (see Session 01)
\item \textbf{VC vs PE vs CVC distinctions} --- know the key differences in stage, control, and return targets
\item \textbf{Share-count based pre/post-money formulas} --- Pre-money $=$ pre-existing shares $\times$ price/share; Post-money $=$ total shares $\times$ price/share
\item \textbf{Statement connections:} retained earnings on balance sheet, CapEx as asset swap (cash $\rightarrow$ fixed asset), cash-flow sign treatment of AR (increase $=$ negative), inventory (increase $=$ negative), AP (increase $=$ positive)
\end{itemize}

\newpage

% ============================================================
% KEY EXAM THEMES
% ============================================================
\section{Key Exam Themes (Cross-Session)}

\begin{enumerate}
\item \textbf{Profit does NOT equal Cash} (Sessions 00, 03, quiz)
\item \textbf{Working Capital formulas and impact on cash flow} (Sessions 00, 03, quiz)
\item \textbf{Operating Cash Flow} = EBIT $-$ Taxes + D\&A $-$ Increase in WC (Session 03, quiz)
\item \textbf{VC Method 5-step process} (Session 04)
\item \textbf{Pre-money / Post-money / Dilution calculations} (Session 04)
\item \textbf{IRR formula and compounding} (Session 04)
\item \textbf{Team $>$ Market $>$ Technology} in investment decisions (Sessions 01, 02)
\item \textbf{Bottom-up financial forecasting}, never top-down (Session 03)
\item \textbf{Business Angel vs VC differences} (Session 01)
\item \textbf{Convertible notes} --- cap, discount, when to use (Session 01)
\item \textbf{Balance sheet must always balance} (Sessions 00, 03)
\item \textbf{Tax credit / loss carry-forward} (Session 03)
\item \textbf{VC fund economics} --- 2\% fee, 20\% carry, hurdle rate, lifecycle (Sessions 01, 03)
\end{enumerate}

\end{document}
